\subsection{A prospective calibration protocol}
\label{sec:decision-rule}

We derive a prospective calibration protocol from the four-model study and
evaluate its decisions descriptively; held-out validation of the thresholds
remains future work. On a new model, the protocol runs a short calibration
packet and measures strict set-leaving, no-gap fraction, static top-$K$
recovery, drift entropy/autocorrelation, and a rotation pilot if available.
It then chooses among static top-$K$, budgeted EMA, rotation, or rejection of
channel-set protection for that regime.

The descriptive decisions match the observed pattern. All four models exceed
0.53 strict set-leaving, so static top-1\% protection needs validation rather
than assumption. Nemotron passes the budgeted-EMA gate: M11b top-10 recovers
0.815 with CI [0.254, 0.923] and beats static top-10 by 0.220 median. The same
gate is ambiguous on DeepSeek and Falcon: top-10 recovery is 0.335 with CI
[-0.409, 0.494] on DeepSeek and 0.044 with CI [-0.144, 0.203] on Falcon. When
a rotation baseline is available, the protocol compares it directly rather
than stacking it with M11b. Granite favors ParoQuant over M11b, and
ParoQuant+M11b is sub-additive.

This protocol is not a universal predictor. Its thresholds are derived from
these four models, so the four-model evaluation is descriptive rather than
held-out validation. The deployment guidance is still concrete: use rotation
where it is locally validated, as on Granite; use budgeted EMA where it clears
a model-local matched-budget gate, as on Nemotron; and avoid claiming a
channel-set positive method when high drift is present but the matched-budget
gate is ambiguous, as on DeepSeek and Falcon. The negative composition result
is part of the guidance: remedies must be selected by measured model behavior,
not stacked by default.
